\documentclass{article}

%**************************************************************
% Importazione package
%**************************************************************

% permette di modificare i margini
%\usepackage[top=3.1cm, bottom=3.1cm, left=2.2cm, right=2.2cm]{geometry}
\usepackage[a4paper]{geometry}

% necessario per risolvere il problema del carattere invisibile per l'a capo
% \DeclareUnicodeCharacter{00A0}{ }

% per scrivere in italiano e in inglese;
% l'ultima lingua (l'italiano) risulta predefinita
\usepackage[english, italian]{babel}

% imposta lo stile italiano per i paragrafi
\usepackage{parskip}

% fornice elenchi numerati in ordine inverso
\usepackage{etaremune}

% comandi per l'appendice
\usepackage{appendix}
\renewcommand\appendixtocname{Appendici}
\renewcommand{\appendixpagename}{Appendici}

% numera anche i paragrafi
\setcounter{secnumdepth}{6}

% elenca anche i paragrafi nell'indice
\setcounter{tocdepth}{6}

% permetti di definire dei colori
\usepackage[usenames,dvipsnames]{color}

% permette di inserire le immagini/tabelle esattamente dove viene usato il
% comando \begin{figure}[H] ... \end{figure}
% evitando che venga spostato in automatico
\usepackage{float}

% tabelle su più pagine
\usepackage{longtable}

% permette l'inserimento di url e di altri tipi di collegamento
\usepackage[colorlinks=true]{hyperref}

\hypersetup{
		breaklinks=true,
    colorlinks=true, % false: boxed links; true: colored links
    citecolor=black,
    filecolor=black,
    linkcolor=black, % color of internal links
    urlcolor=Maroon  % color of external links
}
\urlstyle{same}

% immagini
\usepackage{graphicx}

% permette di riferirsi all'ultima pagina con \pageref{LastPage}
\usepackage{lastpage}

% per avere dei comandi in più da poter usare sulle tabelle
\usepackage{booktabs}

% tabelle con il campo X per riempire lo spazio rimanente sulla riga
\usepackage{tabularx}

% multirow per tabelle
\usepackage{multirow}

% colore di sfondo per le celle
\usepackage[table]{xcolor}
%\rowcolors{2}{gray!25}{} % Colora righe alterne in grigio (non funziona bene)

% permette di fare longtable larghe tutta la pagina (parametro x)
% su Ubuntu non si può installare il pacchetto, deve essere in model/
\usepackage{tabu}

% imposta lo spazio tra le righe di una tabella
\setlength{\tabulinesep}{6pt}

% definisci un nuovo tipo di colonna P che permette di andare a capo con \newline
% e giustificata a sinistra
\usepackage{array}
\usepackage{ragged2e}
\newcolumntype{P}[1]{>{\centering \hspace{0pt}}m{#1}}

% personalizza l'intestazione e piè di pagina
\usepackage{fancyhdr}

% permette di inserire caratteri speciali
\usepackage{textcomp}

% permette di aggiustare i margini e centrare tabelle e figure
\usepackage{changepage}

%Permette di includere i grafici a barre
%IMPORTANTE: deve essere caricato prima di /pgfgantt altrimenti causa conflitto
\usepackage{pgfplots}

% permette di includere i diagrammi Gantt
% su Ubuntu non si può installare il pacchetto, deve essere in model/
\usepackage{pgfgantt}

% permette di includere i grafici a torta
\usepackage{pgf-pie}

% necessario per pgf-pie
\usepackage{tikz}

% permette i path delle immagini con gli spazi
\usepackage{grffile}

% ruota le immagini
\usepackage{rotating}

% permetti di calcolare le larghezze facendo calcoli
\usepackage{calc}


%%
%% Definizione di path globali
%%
% Path assoluta a questa directory, va passato come argomento il resto della path
\newcommand{\ModelPath}[1]{/data/sources/model/#1}

% Path per accedere a modelassets
\newcommand{\ModelAssets}[1]{\ModelPath{modelassets}/#1}

% Path globali per la ricerca di immagini
\graphicspath{
	{/data/logo/PNG}
	{/data/logo/Square}
}



\fancypagestyle{plain}{
	% cancella tutti i campi di intestazione e piè di pagina
	\fancyhf{}
	\lhead{
		\includegraphics[width=1.5cm, keepaspectratio=true]{argo_icona.png}
		\parbox[b]{5cm}{
			\emph{\GroupName{}} \vspace{0pt} \\
			\emph{Progetto \ProjectName{}} \vspace{0pt}
		}
	}
	\chead{}


	\lfoot{
		\DocTitle{} \\
		% differenzia a seconda che \DocVersion{} stampi testo o no
		\setbox0=\hbox{\DocVersion{}\unskip}\ifdim\wd0=0pt
			% nulla
		\else
			\faIcon{tag} v \DocVersion{}
		\fi
	}
	\rfoot{\thepage{} di \pageref{LastPage}}

	% Visualizza una linea orizzontale in cima e in fondo alla pagina
	\renewcommand{\headrulewidth}{0.3pt}
	\renewcommand{\footrulewidth}{0.3pt}
}
\setlength{\headheight}{30pt}
\pagestyle{plain}

% allarga l'header a tutta la pagina
%\fancyhfoffset[L]{\oddsidemargin + \hoffset + 1in}
%\fancyhfoffset[R]{\evensidemargin + \marginparwidth - \marginparsep}

% Per inserire del codice sorgente formattato
\usepackage{listings}
\definecolor{darkgray}{rgb}{.4,.4,.4}
\definecolor{purple}{rgb}{0.65, 0.12, 0.82}
\definecolor{mygreen}{rgb}{0,0.6,0}
\definecolor{mygray}{rgb}{0.5,0.5,0.5}
\definecolor{mymauve}{rgb}{0.58,0,0.82}

\lstset{
  extendedchars=true,          % lets you use non-ASCII characters
  inputencoding=utf8,   % converte i caratteri utf8 in latin1, richiede \usepackage{listingsutf8} anzichè listings
  basicstyle=\ttfamily,        % the size of the fonts that are used for the code
  breakatwhitespace=false,     % sets if automatic breaks should only happen at whitespace
  breaklines=true,             % sets automatic line breaking
  captionpos=b,                % sets the caption-position to top
  commentstyle=\color{mygreen},   % comment style
  frame=none,               % adds a frame around the code
  keepspaces=true,            % keeps spaces in text, useful for keeping indentation of code (possibly needs columns=flexible)
  keywordstyle=\color{blue}\bfseries,     % keyword style
  numbers=none,               % where to put the line-numbers; possible values are (none, left, right)
  numbersep=5pt,              % how far the line-numbers are from the code
  numberstyle=\color{mygray}, % the style that is used for the line-numbers
  rulecolor=\color{black},    % if not set, the frame-color may be changed on line-breaks within not-black text (e.g. comments (green here))
  showspaces=false,           % show spaces everywhere adding particular underscores; it overrides 'showstringspaces'
  showstringspaces=false,     % underline spaces within strings only
  showtabs=false,             % show tabs within strings adding particular underscores
  stepnumber=5,               % the step between two line-numbers. If it's 1, each line will be numbered
  stringstyle=\color{red},    % string literal style
  tabsize=4,                  % sets default tabsize
  firstnumber=1      % visualizza i numeri dalla prima linea
}

% Permetti di utilizzare il grassetto per i caratteri Typewriter (per es. il font di \code{...} e \file{...})
\usepackage[T1]{fontenc}
\usepackage{lmodern}

% Permette di resettare il contatore delle note a piè di pagina ad ogni pagina
\usepackage[perpage, bottom, stable]{footmisc}
\patchcmd{\footref}{\ref}{\ref*}{}{}

% Per poter usare il carattere ° (anche nelle formule matematiche)
\usepackage{gensymb}
\usepackage{newunicodechar}
\newunicodechar{°}{\degree}

% Per importare svg
\usepackage{svg}
% \setsvg{inkscape = inkscape -z -D}

% Per numerare le tabelle e le figure con la sezione in cui si trovano
\usepackage{amsmath}
\numberwithin{figure}{section}
\numberwithin{table}{section}

% Per aggiungere un nuovo livello di profondita (subsubparagraph)

\makeatletter
\newcounter{subsubparagraph}[subparagraph]
\renewcommand\thesubsubparagraph{%
	\thesubparagraph.\@arabic\c@subsubparagraph}
\newcommand\subsubparagraph{%
	\@startsection{subsubparagraph}    % counter
	{6}                              % level
	{\parindent}                     % indent
	{3.25ex \@plus 1ex \@minus .2ex} % beforeskip
	{-1em}                           % afterskip
	{\normalfont\normalsize\bfseries}}
\newcommand\l@subsubparagraph{\@dottedtocline{6}{12em}{5em}}
\newcommand{\subsubparagraphmark}[1]{}
\def\endsubsubparagraph{\par
	\if@noskipsec
	\ifx\@svsechd\@undefined\else\leavevmode\fi
	\fi}
\makeatletter
\def\toclevel@subsubparagraph{6}
\makeatother
% Use \ul{arg} to undlerline text
\usepackage{soul}

% Per le pagine in orientamento orizzontale
\usepackage{pdflscape}
\usepackage{afterpage}

% Set vertical space in tables
\def\arraystretch{1.5}

% Modifica font
\usepackage{fontspec}

\setmainfont{Poppins}[
    Path=\ModelAssets{poppins/} ,
    Extension = .otf ,
    UprightFont=*-Regular ,
    BoldFont=*-Bold ,
    ItalicFont=*-Italic ,
    BoldItalicFont=*-BoldItalic
]

% Aggiunta di icone nel testo dalla libreria fontawesome
\usepackage{fontawesome5}

% Personalizzazione della formattazione dei titoli
\usepackage{titlesec}
\usepackage{titletoc}

% Rappresentazione grafica delle strutture di directory e file
\usepackage{dirtree}